%=======================================================================
%  CSE558 Data Science - Midterm Project Report
%  Monsoon 2026, IIIT-Delhi
%  Built from the course template. Every number below is produced by the
%  code in src/ (results/*.json); see the README.
%  Compile with pdfLaTeX (or Tectonic), twice.
%=======================================================================
\documentclass[10pt,a4paper,twocolumn]{article}

%========================= EDIT THIS =========================
\newcommand{\projecttitle}{Follow the money: how fraud-linked Bitcoin
  transactions and wallets behave differently}
\newif\ifshowhints
\showhintsfalse     % \showhintstrue shows the template's grey instruction lines
%=============================================================

%-------------------------------------------- fonts (no need to edit)
\usepackage{iftex}
\ifPDFTeX
  \usepackage[T1]{fontenc}
  \usepackage[utf8]{inputenc}
  \IfFileExists{carlito.sty}%
    {\usepackage[sfdefault]{carlito}}%              Calibri look
    {\usepackage[scaled=0.92]{helvet}\renewcommand{\familydefault}{\sfdefault}}%
  \usepackage{courier}
\else
  \usepackage[no-math]{fontspec}
  \setmainfont{Carlito}[Extension=.ttf, UprightFont=*-Regular, BoldFont=*-Bold,
    ItalicFont=*-Italic, BoldItalicFont=*-BoldItalic]
\fi

%----------------------------------------------------------- packages
\usepackage[margin=18mm,bottom=20mm,columnsep=7mm]{geometry}
\usepackage{amsmath}
\usepackage{xcolor}
\usepackage{graphicx}
\usepackage{tabularx}
\usepackage{array}
\usepackage{booktabs}
\usepackage{enumitem}
\usepackage{titlesec}
\usepackage{needspace}
\usepackage[font=small,labelfont=bf]{caption}
\usepackage[hidelinks]{hyperref}
\hypersetup{pdftitle={\projecttitle},
            pdfsubject={CSE558 Data Science, Monsoon 2026: Midterm Project Report}}

%------------------------------------------------------------- layout
\definecolor{eg}{HTML}{7396C4}         % worked example (light blue)
\definecolor{hintgrey}{gray}{0.58}     % instructions (light grey)
\definecolor{todo}{HTML}{C0392B}       % fields still to fill in

\titleformat{\section}{\normalfont\large\bfseries}{\thesection.}{0.5em}{}
\titlespacing*{\section}{0pt}{12pt}{3pt}
\titleformat{\subsection}{\normalfont\normalsize\bfseries}{\thesubsection}{0.5em}{}
\titlespacing*{\subsection}{0pt}{10pt}{2pt}
\newcommand{\sectionbreak}{\Needspace{10\baselineskip}}
\newcommand{\subsectionbreak}{\Needspace{10\baselineskip}}

\setlength{\parindent}{0pt}
\setlength{\parskip}{4pt}
\setlist{leftmargin=1.2em, itemsep=2pt, topsep=2pt, parsep=0pt}
\renewcommand{\arraystretch}{1.25}
\setlength{\tabcolsep}{4pt}
\newcolumntype{L}[1]{>{\raggedright\arraybackslash}p{#1}}
\newcolumntype{Y}{>{\raggedright\arraybackslash}X}
\newcolumntype{K}[1]{>{\raggedright\arraybackslash\bfseries\boldmath}p{#1}}  % bold labels

%--------------------------------------------------------- shorthands
\newcommand{\hint}[1]{\ifshowhints{\footnotesize\itshape\color{hintgrey}#1\par}\fi}
\newcommand{\eg}[1]{\textcolor{eg}{#1}}
\newcommand{\egblock}[1]{{\color{eg}#1\par}}
\newcommand{\field}[2]{\textbf{#1}~#2\par}
\newcommand{\tpart}[2]{\par\Needspace{7\baselineskip}\smallskip\textbf{#1)~#2}\par}
\newcommand{\minihead}[1]{\par\Needspace{8\baselineskip}\smallskip\textbf{#1}\par}
\newcommand{\code}[1]{\texttt{#1}}
\newcommand{\todo}[1]{\textcolor{todo}{\textbf{TODO:} #1}}
\newcommand{\figfile}[2]{%
  \IfFileExists{figs/#1}%
    {\includegraphics[width=\linewidth,height=#2,keepaspectratio]{#1}}%
    {\fbox{\parbox[c][#2][c]{\dimexpr\linewidth-2\fboxsep-2\fboxrule\relax}%
      {\centering\small\itshape\color{hintgrey}Your plot: \detokenize{#1}}}}}
\graphicspath{{figs/}}

%=======================================================================
\begin{document}

\twocolumn[%
{\small\bfseries CSE558 Data Science\ \ \textbullet\ \ Midterm Project Report\ \ \textbullet\ \ Monsoon 2026\par}
\vspace{4pt}
{\LARGE\bfseries\raggedright\projecttitle\par}
\vspace{6pt}
\rule{\linewidth}{1pt}

{\small
\begin{tabularx}{\linewidth}{@{}K{2.4cm} Y@{}}
\toprule
Group name & Innovate and Conquer \\
Dataset & Elliptic++ (Elmougy \& Liu, KDD 2023): \url{github.com/git-disl/EllipticPlusPlus} \\
Code & \url{https://github.com/AnantJainn/Midterm-Project} (scripts \code{src/01--03}, README, results) \\
\bottomrule
\end{tabularx}}

\vspace{2pt}
{\small
\begin{tabularx}{\linewidth}{@{}L{0.5cm} L{3.4cm} L{2.3cm} L{4.8cm} Y@{}}
\toprule
\textbf{\#} & \textbf{Name} & \textbf{Roll no.} & \textbf{Institute email} & \textbf{Contribution} \\
\midrule
1 & Anant Jain & PhD25002 & anantj@iiitd.ac.in & Dataset and size check, data preparation and before/after (1, 2, 3); Tests 1, 2 and 4 (5); code and README (6) \\
2 & Gauranshi Gupta & PhD25005 & gauranshig@iiitd.ac.in & EDA, basic questions and problem statements (4); Tests 3 and 5 (5); limitations; report writing \\
\bottomrule
\end{tabularx}}
\vspace{12pt}
]

%-----------------------------------------------------------------------
\section{Dataset}

\field{Source:}{Elliptic++, released by Georgia Tech with the KDD 2023 paper of Elmougy and
Liu, which extends the Elliptic Bitcoin dataset of Weber et al.\ (2019). Downloaded on 21
October 2026 from the Google Drive folder linked in the repository above. Each row is a real
Bitcoin transaction or wallet address; the blockchain analytics firm Elliptic labelled
entities as \emph{illicit} (scams, ransomware, malware, Ponzi schemes, dark markets),
\emph{licit} (exchanges, wallet providers, miners) or \emph{unknown}. The data covers 49 time
steps of about two weeks each.}

\field{What is raw about it:}{the labels sit in separate files from the features, coded
1/2/3; 77\% of rows are \emph{unknown}; 965 transactions miss all 17 named features; the
wallet file repeats an address once for every time step it appears in; amounts and fees are
extremely skewed; the money trail is only an edge list (who paid whom).}

{\small
\begin{tabularx}{\linewidth}{@{}L{2.6cm} Y Y@{}}
\toprule
 & \textbf{Raw file} & \textbf{After encoding} \\
\midrule
\multicolumn{3}{@{}l}{\textit{Transactions} (\code{txs\_features} + \code{txs\_classes})} \\
Instances (rows) & 203{,}769 & 202{,}804 \\
Attributes (columns) & 185 & 196 \\
Instances $\times$ attributes & 37{,}697{,}265 & 39{,}749{,}584 \\
\midrule
\multicolumn{3}{@{}l}{\textit{Wallets} (\code{wallets\_features} + \code{wallets\_classes})} \\
Instances (rows) & 1{,}268{,}260 & 822{,}942 \\
Attributes (columns) & 58 & 66 \\
Instances $\times$ attributes & 73{,}559{,}080 & 54{,}314{,}172 \\
\bottomrule
\end{tabularx}}

\field{Size check:}{transactions: $202{,}804\times196=39{,}749{,}584\ge3{,}000{,}000$;
wallets: $822{,}942\times66=54{,}314{,}172\ge3{,}000{,}000$. Both tables meet the requirement
on their own. Even the 17 named transaction features alone give
$202{,}804\times17=3{,}447{,}668$.}

\minihead{Labels after cleaning}
{\footnotesize\setlength{\tabcolsep}{3pt}
\begin{tabularx}{\linewidth}{@{}L{1.5cm} Y Y Y@{}}
\toprule
 & \textbf{Illicit} & \textbf{Licit} & \textbf{Unknown} \\
\midrule
Transactions & 4{,}545 (2.2\%) & 41{,}500 (20.5\%) & 156{,}759 (77.3\%) \\
Wallets & 14{,}266 (1.7\%) & 251{,}088 (30.5\%) & 557{,}588 (67.8\%) \\
\bottomrule
\end{tabularx}}

Unknown rows stay in the data but not in the tests (Section~2). The labelled rows that the
tests use still meet the size requirement: $46{,}045\times196=9{,}024{,}820$ transaction cells
and $265{,}354\times66=17{,}513{,}364$ wallet cells.

\minihead{Key features}
{\footnotesize
\begin{tabularx}{\linewidth}{@{}L{2.6cm} L{1.3cm} Y@{}}
\toprule
\textbf{Feature} & \textbf{Type} & \textbf{Real-world meaning} \\
\midrule
\code{class} & nominal & illicit / licit / unknown; our response \\
\code{num\_output\_}\allowbreak\code{addresses} & discrete & how many addresses receive money in the transaction \\
\code{num\_input\_}\allowbreak\code{addresses} & discrete & how many addresses pay into it \\
\code{fees}, \code{size} & continuous & fee paid to the miner (BTC); transaction size (bytes) \\
\code{fee\_rate} (derived) & continuous & \code{fees}/\code{size} in satoshi per byte: the price paid for block space \\
\code{total\_BTC} & continuous & total value moved (BTC) \\
\code{Time step} & discrete & 1--49, about two weeks each \\
\code{num\_txs\_as\_sender}, \code{\_receiver} & discrete & a wallet's transactions sending / receiving money \\
\code{lifetime\_}\allowbreak\code{in\_blocks} & discrete & blocks between a wallet's first and last appearance \\
\bottomrule
\end{tabularx}}

%-----------------------------------------------------------------------
\section{Data preparation}

\begin{table*}[!t]
\caption{Data preparation: every issue we found in the raw files, what we did, and why.}
\label{tab:prep}
{\footnotesize
\begin{tabularx}{\textwidth}{@{}L{3.6cm} L{5.4cm} Y@{}}
\toprule
\textbf{Issue found} & \textbf{What we did} & \textbf{Why, and not the alternative} \\
\midrule
Labels in separate files, coded 1/2/3
 & joined on \code{txId} / \code{address} (one-to-one, checked); mapped to names; one-hot
   encoded (\code{is\_illicit}, \code{is\_licit}, \code{is\_unknown})
 & a code like 3 is not a number: treating it as one would make ``unknown'' three times
   ``illicit''. \\
\emph{Unknown} class (77\% of transactions)
 & kept in the file, left out of every test
 & unknown is not ``licit'': counting it as licit would dilute the fraud signal with
   unlabelled fraud. \\
965 transactions miss all 17 named features
 & dropped (519 licit, 446 unknown, no illicit)
 & the 17 are missing \emph{together}, a structural gap in the Elliptic++ augmentation, so
   imputing would invent whole transactions; no illicit row is lost. \\
Wallet file: 347{,}569 exact duplicate rows
 & dropped
 & they add no information and would count a wallet twice in every test. \\
61{,}487 wallets still listed in several time steps
 & collapsed to one row per address; first time step kept
 & we checked that every wallet feature is constant across these rows: they are whole-life
   summaries, so the repeats are the same wallet counted again. \\
Fee rate heavy-tailed (11{,}303 above the IQR fence, max 22{,}614 sat/B)
 & kept and flagged; tests on fee rate use $\log_{10}$
 & these are real, urgent payments; capping them would erase exactly the spread that Test 3
   measures. \\
67 transactions with zero fee
 & kept; left out only where $\log_{10}$ fee rate is needed
 & a zero fee is real (miners accepted free transactions), but has no logarithm. \\
Money trail only an edge list
 & derived \code{receives\_}\allowbreak\code{from\_}\allowbreak\code{illicit}: the transaction spends an output of an illicit
   transaction
 & turns the graph into a feature that a test can use (Test 4). \\
\bottomrule
\end{tabularx}}
\end{table*}

Table~\ref{tab:prep} lists each issue we found in the raw files, what we did about it, and
why we chose that over the alternative.

We also checked Bitcoin's accounting identity, inputs $=$ outputs $+$ fee: it holds for 100\%
of rows, so amounts were not corrupted. There were no duplicate transactions.

%-----------------------------------------------------------------------
\section{What changed after preparation}

{\footnotesize
\begin{tabularx}{\linewidth}{@{}L{3.0cm} L{1.7cm} Y@{}}
\toprule
\textbf{Measure} & \textbf{Before} & \textbf{After} \\
\midrule
Transaction rows & 203{,}769 & 202{,}804 (965 dropped) \\
Transaction attributes & 185 & 196 (11 derived or encoded) \\
Missing cells (named features) & 16{,}405 & 0 \\
Wallet rows & 1{,}268{,}260 & 822{,}942 (one per address) \\
Wallet duplicate rows & 347{,}569 & 0 \\
Wallets counted more than once & 61{,}487 & 0 \\
Wallet attributes & 58 & 66 (pattern and one-hot class) \\
Fee-rate outliers & 11{,}303 & 11{,}303, flagged and analysed on $\log_{10}$ \\
Labelled transactions for tests & 46{,}564 & 46{,}045 (4{,}545 illicit, 41{,}500 licit) \\
Labelled wallets for tests & 265{,}354 & 265{,}354 (14{,}266 illicit, 251{,}088 licit) \\
\bottomrule
\end{tabularx}}

The wallet table shrank by 35\%, and no address was lost: all 822{,}942 remain, each counted
once. On the log scale the fee rate is close to symmetric (skewness 0.27 for licit
transactions), whereas the raw amounts span eight orders of magnitude
($2.7\times10^{-5}$ to 6{,}335 BTC). The encoded label and
the new money-trail and fee-rate features are what make Tests 3 and 4 possible.

%-----------------------------------------------------------------------
\section{Basic questions and problem statement}

\smallskip\textbf{Basic questions}\par
\begin{enumerate}[label=\arabic*.]
  \item \emph{How common is fraud?} 9.9\% of labelled transactions are illicit (4{,}545 of
    46{,}045); 77\% of all transactions carry no label.
  \item \emph{Does fraud move more money?} No: the median illicit transaction moves 0.12 BTC,
    the median licit one 0.48 BTC.
  \item \emph{Is fraud steady over time?} No: the illicit share of a time step ranges from
    0.3\% (step 46) to 36.5\% (step 13) (Figure~\ref{fig:time}).
  \item \emph{Does fraud pay more in fees?} Per byte slightly more (median 105 vs 90 sat/B),
    but illicit transactions are smaller (median 192 vs 372 bytes), so the median fee is
    lower (0.00023 vs 0.0004 BTC).
\end{enumerate}

\begin{figure}[!ht]
  \figfile{fig_timeline.pdf}{4.3cm}
  \caption{The illicit share swings from step to step and falls to near zero right after the
  dark-market shutdown at step 43 reported by Weber et al.; the shaded windows are those of
  Test 5.}
  \label{fig:time}
\end{figure}

\begin{figure}[!ht]
  \figfile{fig_outputs.pdf}{4.0cm}
  \caption{65\% of illicit transactions pay a single receiver; 76\% of licit ones pay two,
  the typical ``payment plus change''.}
  \label{fig:out}
\end{figure}

\begin{figure}[!ht]
  \figfile{fig_feerate.pdf}{3.4cm}
  \caption{Illicit fee rates pile up in a narrow spike at about 100--120 sat/B; licit fee
  rates spread over two orders of magnitude.}
  \label{fig:fee}
\end{figure}

\begin{figure}[!ht]
  \figfile{fig_wallets.pdf}{3.4cm}
  \caption{63\% of illicit wallets both receive and send money (pass-through), against 11\%
  of licit wallets, most of which only receive.}
  \label{fig:wal}
\end{figure}

\Needspace{13\baselineskip}
\minihead{Our problem statements}
\noindent\fbox{\begin{minipage}{\dimexpr\linewidth-2\fboxsep-2\fboxrule\relax}
\textbf{P1 (primary).} Does illicit money move in a different \emph{shape} from legitimate
money: forwarded whole to a single receiver, through wallets that pass it straight on?
(Tests 1, 2)

\smallskip
\textbf{P2.} Does fraud leave a \emph{fingerprint} in the fees: are illicit fee rates more
uniform, as automated software with a fixed fee would produce? (Test 3)

\smallskip
\textbf{P3.} Does fraud \emph{follow the money trail}: is a transaction that spends illicit
money itself more likely to be illicit, comparing like with like in time? (Test 4)

\smallskip
\textbf{P4.} Did the dark-market shutdown at step 43 lower the illicit share? (Test 5)
\end{minipage}}

\field{How they are novel:}{the basic questions compare totals and averages, one row at a
time. Ours ask about \emph{behaviour}: how money is split (outputs), how a wallet uses it
(receive / send), how \emph{variable} the fees are rather than how high, and what the
\emph{previous} transaction on the money trail was, which no single row records and has to
be derived from the graph. Elliptic++ has so far been used mainly to train classifiers; we
test these behaviours directly, with effect sizes.}

%-----------------------------------------------------------------------
\section{Statistical inferences}

\field{Rows tested:}{only labelled rows: 46{,}045 transactions (4{,}545 illicit, 41{,}500
licit) and 265{,}354 wallets (14{,}266 illicit, 251{,}088 licit). All tests use
$\alpha=0.05$. Hypotheses, sides and the smallest differences that matter were fixed in the
code before any test statistic was computed. Where a p-value is below the smallest number a
computer can store ($\approx10^{-308}$), we give its order of magnitude from the log
p-value.}

%------------------------------------------------------------ Test 1

\Needspace{12\baselineskip}
\subsection*{Test 1 (P1): do illicit transactions pay a single receiver more often?}

\tpart{a}{Hypotheses, written before running the test}
{\small
\begin{tabularx}{\linewidth}{@{}K{3.4cm} Y@{}}
\toprule
$H_0$ & $p_{\text{ill}}-p_{\text{lic}}\le0$ \\
$H_1$ & $p_{\text{ill}}-p_{\text{lic}}>0$ \\
Significance level & $\alpha=0.05$ \\
Side & one-sided (right) \\
Test & two-proportion $z$-test (pooled) \\
Smallest difference that matters & 10 percentage points \\
\bottomrule
\end{tabularx}}

$p$: share of transactions with exactly one output address. One-sided, because P1 claims
illicit money is forwarded whole. A transaction with one receiver and no change output is
the signature of ``sweeping'' an address empty; 10 points is the gap at which an analyst
screening transactions would gain from the rule.

\tpart{b}{Other tests we considered}
The response is binary, the two groups are independent and large (the smallest expected cell
count is 564), so the normal approximation of the $z$-test holds, and it gives the difference
in proportions directly. With $x$ single-output transactions out of $n$ in each class,
$\hat p=x/n$:
{\small
\[ z=\frac{\hat p_{\text{ill}}-\hat p_{\text{lic}}}
   {\sqrt{\bar p(1-\bar p)\bigl(\tfrac{1}{n_{\text{ill}}}+\tfrac{1}{n_{\text{lic}}}\bigr)}},
   \qquad \bar p=\frac{x_{\text{ill}}+x_{\text{lic}}}{n_{\text{ill}}+n_{\text{lic}}}, \]}%
$p=1-\Phi(z)$. The Newcombe interval for $d=\hat p_{\text{ill}}-\hat p_{\text{lic}}$ combines
the Wilson score intervals $(l_i,u_i)$ of the two proportions:
{\small
\[ L=d-\sqrt{(\hat p_{\text{ill}}-l_{\text{ill}})^2+(u_{\text{lic}}-\hat p_{\text{lic}})^2}, \]
\[ U=d+\sqrt{(u_{\text{ill}}-\hat p_{\text{ill}})^2+(\hat p_{\text{lic}}-l_{\text{lic}})^2}. \]}%
\par
{\footnotesize
\begin{tabularx}{\linewidth}{@{}L{2.3cm} Y@{}}
\toprule
\textbf{Also considered} & \textbf{Why we did not use it} \\
\midrule
Pearson $\chi^2$ (2$\times$2) & the same test squared ($\chi^2=12{,}780=z^2$, d.f.~1), but
two-sided only. \\
Fisher's exact test & exact, but built for small counts; here it agrees (OR $=25.9$,
$p<10^{-300}$) and adds nothing. \\
Welch $t$ on \code{num\_output\_}\allowbreak\code{addresses} & compares means of a count driven by licit
batch payouts with up to 13{,}107 outputs (mean 6.17 vs 1.35); it answers a different
question. \\
\bottomrule
\end{tabularx}}

\tpart{c}{Result}
{\small
\begin{tabularx}{\linewidth}{@{}K{3.4cm} Y@{}}
\toprule
Estimate & $+58.2$ points (64.9\% illicit, 6.7\% licit) \\
95\% confidence interval (Newcombe) & $+56.8$ to $+59.6$ points \\
Test statistic (d.f.) & $z=113.0$ (none: normal reference) \\
p-value & $\approx10^{-2778}$ \\
Decision & reject $H_0$ \\
\bottomrule
\end{tabularx}}

Illicit transactions pay a single receiver nearly ten times as often as licit ones. The whole
interval lies far beyond 10 points, so the difference matters. Together with the 89\% of
illicit transactions that have a single input (licit: 50\%), the typical illicit transaction
is ``one address in, one address out''.

%------------------------------------------------------------ Test 2

\Needspace{12\baselineskip}
\subsection*{Test 2 (P1): do illicit wallets pass money straight through?}

\tpart{a}{Hypotheses, written before running the test}
{\small
\begin{tabularx}{\linewidth}{@{}K{3.4cm} Y@{}}
\toprule
$H_0$ & $\pi_{\text{ill},k}=\pi_{\text{lic},k}$ for every pattern $k$ (independence) \\
$H_1$ & $\pi_{\text{ill},k}\ne\pi_{\text{lic},k}$ for at least one $k$ \\
Significance level & $\alpha=0.05$ \\
Side & upper tail of $\chi^2$ (any departure) \\
Test & Pearson $\chi^2$ test of independence, 2$\times$3 \\
Smallest difference that matters & Cram\'er's $V\ge0.10$ and a 10-point gap in the
pass-through share \\
\bottomrule
\end{tabularx}}

$\pi_{c,k}$: share of class-$c$ wallets with pattern $k$: \emph{receive only}, \emph{pass-through} (receives and sends) or \emph{send only}
(received before the data window). $V=0.10$ is the conventional ``small'' association; we
add a 10-point gap so that the effect can be read in wallets.

\tpart{b}{Other tests we considered}
Two nominal variables, one with three levels, and independent wallets: this is the setting
of the $\chi^2$ test of independence. With observed counts $O_{ik}$ (class $i$, pattern $k$),
row totals $R_i$, column totals $C_k$ and $n$ wallets:
{\small
\[ \chi^2=\sum_{i=1}^{2}\sum_{k=1}^{3}\frac{(O_{ik}-E_{ik})^2}{E_{ik}},\qquad
   E_{ik}=\frac{R_iC_k}{n}, \]
\[ \nu=(2-1)(3-1)=2,\qquad V=\sqrt{\frac{\chi^2}{n\,(\min(2,3)-1)}}, \]
\[ r_{ik}=\frac{O_{ik}-E_{ik}}{\sqrt{E_{ik}}}. \]}%
$r_{ik}$ is the Pearson residual of each cell. The smallest expected count is 2{,}035, far
above 5.
\par
{\footnotesize
\begin{tabularx}{\linewidth}{@{}L{2.3cm} Y@{}}
\toprule
\textbf{Also considered} & \textbf{Why we did not use it} \\
\midrule
$G$-test (likelihood ratio) & asymptotically equivalent; same decision ($G=21{,}428$,
d.f.~2). Pearson's is the standard to report. \\
Two-proportion $z$ on pass-through only & throws away the other two patterns; it agrees
($z=172.1$). We use it for the estimate below. \\
Mann--Whitney on \code{total\_txs} & tests whether one class has more transactions, not
\emph{how} wallets use money. \\
\bottomrule
\end{tabularx}}

\tpart{c}{Result}
{\small
\begin{tabularx}{\linewidth}{@{}K{3.4cm} Y@{}}
\toprule
Estimate & pass-through share $+51.8$ points (63.3\% vs 11.5\%); Cram\'er's $V=0.338$ \\
95\% confidence interval & $+51.0$ to $+52.6$ points; $V$: 0.334 to 0.343 (bootstrap) \\
Test statistic (d.f.) & $\chi^2=30{,}397$ (2) \\
p-value & $\approx10^{-6601}$ \\
Decision & reject $H_0$ \\
\bottomrule
\end{tabularx}}

The association is moderate ($V\approx0.34$) and both intervals clear our thresholds. The
Pearson residuals $r_{ik}$ show where it comes from: illicit pass-through $+155$, illicit
receive-only $-67$. Fraud-linked wallets act as relays, not as savings.

%------------------------------------------------------------ Test 3

\Needspace{12\baselineskip}
\subsection*{Test 3 (P2): are illicit fee rates less variable?}

\tpart{a}{Hypotheses, written before running the test}
{\small
\begin{tabularx}{\linewidth}{@{}K{3.4cm} Y@{}}
\toprule
$H_0$ & $\sigma_{\text{ill}}\ge\sigma_{\text{lic}}$ \\
$H_1$ & $\sigma_{\text{ill}}<\sigma_{\text{lic}}$ \\
Significance level & $\alpha=0.05$ \\
Side & one-sided (left) \\
Test & Brown--Forsythe test \\
Smallest difference that matters & SD ratio $\sigma_{\text{ill}}/\sigma_{\text{lic}}\le0.80$ \\
\bottomrule
\end{tabularx}}

$\sigma$: standard deviation of $\log_{10}$ fee rate (4{,}545 illicit, 41{,}452 licit; 48
zero-fee licit rows excluded). One-sided, because P2 claims uniformity. A spread 20\% narrower
is what a monitoring rule could pick up.

\tpart{b}{Other tests we considered}
We compare spreads, not means. Brown--Forsythe (Levene with the median) runs a $t$-test on
the absolute deviations from each group's median, so it does not need normality and gives a
one-sided p-value. For log fee rate $x_{ij}$ of transaction $i$ in group $j$, with group
median $\tilde x_j$:
{\small
\[ z_{ij}=|x_{ij}-\tilde x_j|,\qquad
   t=\frac{\bar z_{\text{ill}}-\bar z_{\text{lic}}}{s_p\sqrt{\tfrac{1}{n_{\text{ill}}}+\tfrac{1}{n_{\text{lic}}}}},
   \qquad \nu=n_{\text{ill}}+n_{\text{lic}}-2, \]}%
where $s_p^2$ is the pooled variance of the $z_{ij}$, and $F=t^2$ is Levene's statistic. The
estimate $s_{\text{ill}}/s_{\text{lic}}$ gets a percentile bootstrap interval: the 2.5th and
97.5th percentiles of the ratio over 2{,}000 resamples drawn within each group.
\par
{\footnotesize
\begin{tabularx}{\linewidth}{@{}L{2.3cm} Y@{}}
\toprule
\textbf{Also considered} & \textbf{Why we did not use it} \\
\midrule
$F$-test of variances & valid only for normal data; illicit log rates are skewed (1.22) and
Shapiro--Wilk rejects normality ($p\approx10^{-59}$). It agrees ($F=0.31$). \\
Bartlett's test & just as sensitive to non-normality ($\chi^2=2{,}051$). \\
Fligner--Killeen & rank-based and also valid; agrees ($\chi^2=2{,}538$), but two-sided only. \\
Wilcoxon signed-rank, by time step & controls for the fee market changing over time; we run
it as a check (below). \\
\bottomrule
\end{tabularx}}

\tpart{c}{Result}
{\small
\begin{tabularx}{\linewidth}{@{}K{3.4cm} Y@{}}
\toprule
Estimate & SD ratio 0.559 (0.255 vs 0.456 in $\log_{10}$) \\
95\% confidence interval & 0.543 to 0.575 (bootstrap, 2{,}000 resamples) \\
Test statistic (d.f.) & $t=-57.87$ (45{,}995); $F=t^2=3{,}349$ \\
p-value & $\approx10^{-729}$ \\
Decision & reject $H_0$ \\
\bottomrule
\end{tabularx}}

Illicit fee rates are about 44\% less spread out, and the whole interval is below 0.80. The
effect is not a by-product of time: within the same time step, illicit fees are the narrower
group in 41 of 45 steps (median SD 0.082 vs 0.206; Wilcoxon $W=106$, $p=2.3\times10^{-7}$).
The fingerprint is concrete: 25.4\% of illicit transactions pay a fee of exactly 0.0002 BTC
(licit: 8.7\%), typically on a 192-byte, one-in-one-out transaction, which gives the spike at
104 sat/B in Figure~\ref{fig:fee}.

%------------------------------------------------------------ Test 4

\Needspace{12\baselineskip}
\subsection*{Test 4 (P3): does fraud follow the money trail?}

\tpart{a}{Hypotheses, written before running the test}
{\small
\begin{tabularx}{\linewidth}{@{}K{3.4cm} Y@{}}
\toprule
$H_0$ & common odds ratio $\text{OR}_{\text{MH}}\le1$ \\
$H_1$ & $\text{OR}_{\text{MH}}>1$ \\
Significance level & $\alpha=0.05$ \\
Side & one-sided (right) \\
Test & Cochran--Mantel--Haenszel (CMH), stratified by time step \\
Smallest difference that matters & OR $\ge2$ \\
\bottomrule
\end{tabularx}}

Exposure: the transaction spends an output of an illicit transaction. Outcome: it is
illicit. We stratify because a busy fraud period (Figure~\ref{fig:time}) makes both the
exposure and the outcome common at once, which alone would create an association. An OR of 2
doubles the odds, the level at which flagging a payer's payees is worth an analyst's time.

\tpart{b}{Other tests we considered}
A 2$\times$2 table in each of 48 time steps, independent transactions: CMH tests a common
odds ratio across strata while comparing like with like in time. In time step $k$, $a_k$ =
exposed illicit, $b_k$ = exposed licit, $c_k$ = unexposed illicit, $d_k$ = unexposed licit,
$n_k=a_k+b_k+c_k+d_k$:
{\small
\[ \chi^2_{\text{CMH}}=\frac{\bigl[\sum_k\bigl(a_k-\mathrm{E}(a_k)\bigr)\bigr]^2}{\sum_k\mathrm{Var}(a_k)},
   \qquad \text{OR}_{\text{MH}}=\frac{\sum_k a_kd_k/n_k}{\sum_k b_kc_k/n_k}, \]
\[ \mathrm{E}(a_k)=\frac{(a_k+b_k)(a_k+c_k)}{n_k}, \]
\[ \mathrm{Var}(a_k)=\frac{(a_k+b_k)(c_k+d_k)(a_k+c_k)(b_k+d_k)}{n_k^2(n_k-1)}. \]}%
$\nu=1$; the one-sided $p$ is $1-\Phi\bigl(\operatorname{sign}(\text{OR}_{\text{MH}}-1)
\sqrt{\chi^2_{\text{CMH}}}\bigr)$. The interval uses the Robins--Breslow--Greenland variance
of $\ln\text{OR}_{\text{MH}}$.
\par
{\footnotesize
\begin{tabularx}{\linewidth}{@{}L{2.3cm} Y@{}}
\toprule
\textbf{Also considered} & \textbf{Why we did not use it} \\
\midrule
Pearson $\chi^2$ / Fisher, pooled & ignore time and so are confounded: crude OR 31.9 vs 19.3
stratified. \\
Logistic regression with step dummies & agrees in direction (OR 24.7, 21.4 to 28.6), but
needs a model with 49 parameters; CMH is the direct test. \\
Breslow--Day (homogeneity) & not a test of $H_0$ but of CMH's assumption; run below. \\
\bottomrule
\end{tabularx}}

\tpart{c}{Result}
{\small
\begin{tabularx}{\linewidth}{@{}K{3.4cm} Y@{}}
\toprule
Estimate & $\text{OR}_{\text{MH}}=19.3$ (73.8\% vs 8.1\% illicit) \\
95\% confidence interval & 16.7 to 22.3 \\
Test statistic (d.f.) & $\chi^2_{\text{CMH}}=2{,}906$ (1) \\
p-value & $\approx10^{-633}$ \\
Decision & reject $H_0$ \\
\bottomrule
\end{tabularx}}

Even within the same two weeks, a transaction that spends illicit money has 19 times the odds
of being illicit itself; risk ratio 9.1 (8.7 to 9.5). Breslow--Day rejects a common OR
($\chi^2=235$, 38 d.f., $p\approx10^{-29}$): the size of the effect varies over time. Its
direction does not: the odds ratio is at least 2 in 46 of the 48 time steps (median 15.3), so
the conclusion stands.

%------------------------------------------------------------ Test 5

\Needspace{12\baselineskip}
\subsection*{Test 5 (P4): did the illicit share fall after the dark-market shutdown?}

\tpart{a}{Hypotheses, written before running the test}
{\small
\begin{tabularx}{\linewidth}{@{}K{3.4cm} Y@{}}
\toprule
$H_0$ & $\mu_{\text{after}}-\mu_{\text{before}}\ge0$ \\
$H_1$ & $\mu_{\text{after}}-\mu_{\text{before}}<0$ \\
Significance level & $\alpha=0.05$ \\
Side & one-sided (left) \\
Test & Welch $t$-test on per-step illicit shares \\
Smallest difference that matters & 5 percentage points \\
\bottomrule
\end{tabularx}}

$\mu$: mean illicit share per time step, steps 36--42 (before) and 43--49 (after), 7 each.
One-sided, because a shutdown should remove fraud. 5 points is about half the pre-shutdown
level, a drop that would show in an exchange's compliance workload.

\tpart{b}{Other tests we considered}
The unit is the time step, not the transaction: transactions in one step share a market, so
pooling them would overstate the evidence. With 7 values per group we need a small-sample
test that does not assume equal variances. With group means $\bar x$, variances $s^2$ and
sizes $n$ ($a$ = after, $b$ = before):
{\small
\[ t=\frac{\bar x_a-\bar x_b}{\sqrt{s_a^2/n_a+s_b^2/n_b}},\qquad
   \nu=\frac{\bigl(s_a^2/n_a+s_b^2/n_b\bigr)^2}
   {\frac{(s_a^2/n_a)^2}{n_a-1}+\frac{(s_b^2/n_b)^2}{n_b-1}}, \]}%
(Welch--Satterthwaite), and the interval is
$(\bar x_a-\bar x_b)\pm t_{0.975,\nu}\sqrt{s_a^2/n_a+s_b^2/n_b}$.
\par
{\footnotesize
\begin{tabularx}{\linewidth}{@{}L{2.3cm} Y@{}}
\toprule
\textbf{Also considered} & \textbf{Why we did not use it} \\
\midrule
Pooled two-proportion $z$ on 15{,}014 transactions & treats transactions as independent:
$z=-15.7$, $p=7\times10^{-56}$. Far too confident. \\
Student $t$ & variances look equal (Levene $p=0.90$), giving the same $t$; Welch was fixed in
advance. \\
Mann--Whitney (exact) & robust to the non-normal ``after'' group (Shapiro $p=0.04$); agrees:
$U=9$, $p=0.027$. \\
Exact permutation test & all 3{,}432 splits; agrees: $p=0.020$. \\
\bottomrule
\end{tabularx}}

\tpart{c}{Result}
{\small
\begin{tabularx}{\linewidth}{@{}K{3.4cm} Y@{}}
\toprule
Estimate & $-5.4$ points (3.7\% after, 9.1\% before) \\
95\% confidence interval & $-10.2$ to $-0.5$ points \\
Test statistic (d.f.) & $t=-2.43$ (11.9) \\
p-value & 0.016 \\
Decision & reject $H_0$ \\
\bottomrule
\end{tabularx}}

The illicit share did fall after step 43, and all four tests agree. Whether the fall
\emph{matters} is not settled: the estimate just passes our 5 points, but the interval
reaches $-0.5$, so a trivial drop cannot be ruled out with 7 steps. The shares at steps 48
and 49 (7.6\%, 11.8\%) show fraud returning, so the effect looks temporary.

%-----------------------------------------------------------------------
\section{Limitations}

\begin{itemize}
  \item \textbf{Labels come from one firm.} ``Illicit'' is Elliptic's judgement, from its own
  heuristics; if those heuristics follow the money trail, Test 4 partly measures the
  labelling. We cannot check this from the data.
  \item \textbf{77\% of transactions are unknown.} Our results hold for labelled transactions,
  which may be the easier cases to label.
  \item \textbf{Association, not cause.} We show how fraud-linked money behaves, not that the
  behaviour causes or proves fraud; licit services (exchanges sweeping deposits) also
  forward money whole.
  \item \textbf{Huge samples.} With 46{,}045 rows almost any difference is significant, which
  is why we judge each result by its interval against a pre-set threshold.
  \item \textbf{Shutdown dating.} Step 43 is the date reported by Weber et al.; we did not
  date it independently, and with 7 steps per side Test 5 has low power.
  \item \textbf{One period.} The data covers 49 two-week steps of Bitcoin; fee markets and
  fraud methods have changed since.
\end{itemize}

%-----------------------------------------------------------------------
\newpage                               % in two-column mode: start the next column
\section{Conclusion}

Fraud-linked Bitcoin money moves differently from legitimate money, and the differences are
large, not merely significant: four of our five tests clear their pre-set thresholds with the
whole 95\% interval. Illicit money is forwarded whole (65\% vs 7\% single-receiver
transactions), through wallets that pass it on (63\% vs 11\%), with a fee fingerprint of
automated software (spread 44\% narrower, a quarter paying exactly 0.0002 BTC), and it
follows the money trail (common odds ratio 19, within the same time step). The fifth test
shows that the dark-market shutdown lowered the illicit share, but we cannot show the fall
was large or lasting.

These are simple, explainable red flags (one address out, pass-through wallet, fixed default
fee, money received from a known illicit transaction) that an analyst can apply without a
trained model. The next step is to check whether they hold for the 156{,}759 unlabelled
transactions and for current Bitcoin data.

{\footnotesize\textbf{References.} M.~Weber et al., ``Anti-Money Laundering in Bitcoin:
Experimenting with Graph Convolutional Networks for Financial Forensics'', KDD Workshop on
Anomaly Detection in Finance, 2019. Y.~Elmougy and L.~Liu, ``Demystifying Fraudulent
Transactions and Illicit Nodes in the Bitcoin Network for Financial Forensics'', KDD 2023.\par}

\end{document}
